\documentclass[11pt]{article}

\usepackage[margin=1in]{geometry}
\usepackage{microtype}
\usepackage{amsmath}
\usepackage{graphicx}
\usepackage{booktabs}
\usepackage[hidelinks]{hyperref}

\graphicspath{{figures/}}

\newcommand{\natex}{\texttt{natex}}

\title{Exact to the Quarter, Lost in the Generation: US AI-Chip Export
Controls in Product-Level Shipment Panels}
\author{Hauke Hillebrandt\thanks{University College London. Email:
\texttt{ucjthhi@ucl.ac.uk}. This paper and the underlying \natex{} software were
prepared with substantial assistance from Anthropic's Claude models; the author
reviewed the analyses and text and is responsible for all remaining errors.}}
\date{July 2026 \\ \small analysis pass of record: 2026-07, \natex{} v0.2.0, seed 0}

\begin{document}

\maketitle

\begin{abstract}
\noindent At product granularity, the US export-control rounds on AI
accelerators localize in time but not in attribution. In chip~$\times$~quarter
panels of shipment flows built from Epoch AI's chip-sales and chip-components
estimates, all six known-treatment difference-in-differences scans put the
maximum-likelihood-ratio break at the declared policy quarter \emph{exactly},
and the October-2023 A800/H800 ban clears significance (Bernoulli scan
$p=0.030$; the banned SKUs' shipments go to literal zero within one quarter),
while the April-2025 H20 licensing round localizes exactly but insignificantly
($p=0.44$--$0.60$ on sales, $0.13$--$0.23$ on components). The H20 effect
magnitude is large and sign-correct --- $\hat\tau=-2.756$ asinh units under
two-way fixed effects, CR1 $t=-4.23$, 0 of 6 matched placebos exceeding, at
the exact placebo-in-space floor $p=0.143$ --- but cannot be causally
attributed: an unsupervised rediscovery scan shows that the panel's sharpest
structure at 2025Q2 is the \emph{entire Hopper-generation wind-down} with the
H20 inside it, and the placebo battery generates same-size effects from
ordinary product-lifecycle churn, with the matched donor pool either
contaminated (3 of 6 controls are Chinese substitutes whose demand the ban
plausibly boosted) or too small (3 clean controls, placebo $p=0.75$). The
date localization of the October-2023 round is the headline positive result;
the correctly identified lifecycle confound and the pipeline's clean refusals
are the rest of the story.
\end{abstract}

\section{Introduction}\label{sec:intro}

On 7 October 2022 the US Bureau of Industry and Security (BIS) imposed the
first broad export controls on advanced AI accelerators destined for China, a
policy contemporaries called a landmark in US--China technology competition
\cite{allen2022}. Nvidia responded within weeks with China-market SKUs --- the
A800 and H800 --- engineered just under the control thresholds. On 17 October
2023 BIS closed that loophole: the updated advanced-computing rule
\cite{bis2023} adjusted the chip-level performance parameters so that the
workaround SKUs themselves became controlled. Nvidia iterated once more with
the H20, a Hopper-generation part compliant with the 2023 thresholds ---
until 9 April 2025, when the US government informed the company that H20
exports to China (and any chip matching its bandwidth profile) would require
a license, indefinitely; Nvidia's securities filing disclosed up to
\$5.5~billion of associated charges \cite{nvidia2025}.

A growing empirical literature evaluates these controls. On the supplier
side, \cite{crosignani2025} show that US firms subject to export controls
halt sales to their Chinese customers as intended but largely fail to replace
them, losing market value, profitability, and employment. On the evasion
side, \cite{fist2023} document underground markets for controlled
accelerators and estimate the 2023 smuggling flow at hundreds to low
thousands of chips. A companion country-level study in this collection
\cite{hillebrandt2026legs} found that the October-2023 round \emph{bent}
China's legal chip-stock growth downward while aggregate compute was
substantially bypassed. What that literature leaves untested is the sharpest
granularity at which the policy is written: the named SKU. Each control round
names specific products; if the rounds bind, product-level shipment panels
should carry a break at the declared quarter in exactly the named SKUs ---
and nowhere else.

This paper runs that test three ways with the \natex{} toolkit \cite{natex}.
Leg~1 asks whether a scan-based difference-in-differences that searches over
break dates \cite{herlands2018}, \emph{knowing} which chips were treated,
recovers the declared policy quarter from the data. Leg~2 removes the
treatment labels entirely and asks what the sharpest structure in each panel
actually is. Leg~3 estimates effect magnitudes with the placebo-based
inference appropriate to a handful of treated clusters
\cite{conleytaber2011,cameron2015}. The answer pattern --- exact localization
everywhere, significance only for October 2023, attribution defeated by
product-lifecycle churn --- is documented below, refusals included.

\section{Data}\label{sec:data}

Two Epoch AI datasets are used. \emph{Data on AI Chip Sales}
\cite{epochchipsales} provides quarterly shipment-flow estimates by chip
model in thousands of H100-equivalents (kH100e); the extract's estimates
were generated in January 2026 with 2025Q4 the last full-coverage quarter.
\emph{Data on AI Chip Components} \cite{epochchipcomponents} provides
quarterly CoWoS advanced-packaging wafer estimates by chip for
2024Q1--2025Q4. From these, three chip~$\times$~quarter panels are built
(time index $t=4\,\mathrm{year}+\mathrm{quarter}-1$, so $8095=$ 2023Q4 and
$8101=$ 2025Q2): (i)~\emph{sales Oct-23}, 7 chips $\times$ 16 quarters
(2022Q1--2025Q4), treated \{A800, H800\}, declared $t_0=$ 2023Q4;
(ii)~\emph{sales H20}, 20 chips $\times$ 8 quarters (2024Q1--2025Q4),
treated \{H20\}, declared $t_0=$ 2025Q2; and (iii)~\emph{components H20}, 15
chips $\times$ 8 quarters of CoWoS wafers, treated \{H20, MI308X\}, declared
$t_0=$ 2025Q2.

The outcome is $\operatorname{asinh}$ of the flow. Zeros are true zeros (no
shipments), not missingness; chips born after $t_0-1$ are excluded at build
time so that treatment cannot be mechanically confounded with birth; and the
ragged 2026 tail is dropped at the 2025Q4 snapshot. Because the treatment
effect here operates on the extensive margin --- banned SKUs' shipments go to
approximately zero --- asinh magnitudes are unit-dependent and should not be
read as percentage effects \cite{chenroth2024}; the qualitative reading
(shipments driven to zero) is the transform-robust statement. One
measurement caveat is carried through the paper: Huawei Ascend and Cambricon
Siyuan flows in the source data are annual allocations smeared uniformly
across quarters, which makes those series smooth by construction.

\section{Design and Methods}\label{sec:methods}

\paragraph{Leg 1: known-treatment scans (date localization).} For each panel,
the \natex{} SuDDDS scan \cite{natex,herlands2018} maximizes a
likelihood-ratio statistic for a declared-treatment DiD over candidate break
quarters $t_0$ and window widths $W\in\{2,4,8\}$. Because the unit column is
the chip itself, no subset search is available to the known-treatment scan:
it localizes over $(t_0,W)$ only, so \emph{date localization is the contest}
--- does the maximum-LLR break land on the declared quarter? The primary
model is Bernoulli with within-cluster calibration (wcc); a normal-model
cross-check is calibrated under a dependence-preserving AR(1)-by-unit null.
Permutation calibration uses $Q=99$ draws at seed 0, so the smallest
attainable scan $p$ is $0.01$.

\paragraph{Leg 2: unsupervised discovery rematch.} The same scan machinery is
then run with the treatment column deleted: $\theta$ is the observed flow, a
categorical chip-name dimension is admitted to the subset search (model
normal, method single\_delta, restarts 8, panel-null $Q=99$, seed 0), and the
scan reports whatever subset~$\times$~window is sharpest. Discovery never
reads treatment labels; the question is whether the policy set is the
panel's sharpest structure or merely sits inside something larger.

\paragraph{Leg 3: effects.} A manual two-way fixed-effects (TWFE) DiD with
chip and quarter fixed effects gives $\hat\tau$ in asinh units with CR1
standard errors clustered by chip. With at most 7 clusters, CR1 $t$-statistics
are known to over-reject \cite{cameron2015} --- and demonstrably do here ---
so the \emph{inference of record} is an exact placebo-in-space test in the
spirit of \cite{conleytaber2011}: every same-size subset of control chips is
pseudo-treated (true treated units dropped) and $p$ is the rank of
$|\hat\tau|$ in that exact distribution. The preferred specification is
lifecycle-matched: donor chips alive (nonzero) in all 8 quarters, removing
births and deaths from the donor pool. Robustness: dropping 2025Q1 (a visible
pull-forward spike of 54.2k H100e before the April license), restricting to
non-Chinese controls (substitution-contamination check), a placebo-in-time
at pseudo-$t_0=$ 2024Q4 on pre-period data only, and an event study relative
to 2025Q1.

\paragraph{Honest refusals, stated as run.} \natex's built-in
dd/synthetic-control/GESS effect estimators refused on all three panels (0
usable placebo units; $\geq 5$ required) --- hence the manual TWFE leg. The
scan's anticipation test refused on all known-treatment runs (no usable
pre-quarters at these window widths; reported as not-passed with
$p=\mathrm{None}$, a refusal rather than a failure). And the
lifecycle-matched Oct-23 effect leg refused outright: 0 of 5 control chips
are alive in all 16 quarters, so no matched donor pool exists. Refusals are
reported as refusals, not patched over.

\section{Results}\label{sec:results}

\begin{figure}[t]
\centering
\includegraphics[width=\textwidth]{fig1}
\caption{(a) Oct-2023 round: $\operatorname{asinh}$ shipment flows for the
banned SKUs (A800, H800) against the five incumbent control chips; the scan
puts the break exactly at the declared quarter (dashed line, 2023Q4 boundary)
with $p=0.030$, and treated shipments hit literal zero within one quarter.
(b) H20 round: the H20 against 6 lifecycle-matched controls (solid) and 13
unmatched chips entering or exiting the market (dotted) --- the churn that
defeats attribution; the break again lands exactly on the declared quarter
but at $p=0.44$. (c) Event study, H20 minus matched controls relative to
2025Q1: flat pre-trend, monotone post-collapse to $-4.07$ asinh units by
2025Q4; the shaded band is the placebo-in-space envelope (max $|\tau|$ over
the 6 matched controls), which the pooled estimate exceeds only at the
$p=1/7$ floor.}
\label{fig:fig1}
\end{figure}

\paragraph{Localization: six for six, one significant.} Every known-treatment
scan puts the maximum-LLR break at the declared quarter exactly
(Table~\ref{tab:main}, panel A). For the Oct-2023 round (7 chips $\times$ 16
quarters, treated A800+H800) the Bernoulli scan discovers $t_0=8095=$ 2023Q4
with $W=4$ and LLR $=4.896$, scan $p=0.030$ ($Q=99$); the normal/ar1\_unit
cross-check finds the same $t_0$ and $W$ with LLR $=4.085$, $p=0.14$. The
banned SKUs' shipments go to literal zero within one quarter
(Figure~\ref{fig:fig1}a). For the H20 round the localization is equally
exact but never significant: sales panel $p=0.44$ (Bernoulli) and $0.60$
(normal); components panel (treated H20+MI308X) $p=0.13$ and $0.23$. The
composition check passes ($p=1.0$) in all six scans, and role-assignment
sanity passes throughout: the scans are dating the declared treatment, not
rediscovering an artifact.

\paragraph{Discovery: the generation, never the policy set.} With treatment
labels removed (Table~\ref{tab:main}, panel B), the Oct-23 panel's sharpest
structure is \{A800, H100/H200, H800, TPU v4, TPU v5e\} at \emph{2022Q4}
(LLR $=16.16$, $p=0.01$ floor) --- the Hopper production ramp and the birth
of the A800, itself the workaround SKU created by the October-\emph{2022}
round. The ban quarter 2023Q4 does not appear in the top five discoveries;
the closest is \{A100, A800, H800\} at 2024Q1 (rank 3) --- the banned SKUs
plus the A100's end-of-life. On the H20 sales panel, discovery puts the
sharpest break at 2025Q2 \emph{exactly} (LLR $=17.32$, $p=0.05$), but the
discovered subset is \{H100/H200, H20, MI300A, TPU v5e, TPU v5p, Trainium1\}
--- the H20 merged into the entire Hopper-generation wind-down, with the
complementary Blackwell-generation upshift as the runner-up. The components
panel repeats the pattern one quarter early (7-chip old-generation set at
2025Q1, LLR $=16.41$, $p=0.02$). Every autonomously discovered subset is a
chip \emph{generation}, never the policy set: at product granularity, the
H20 ban and the Hopper-to-Blackwell transition are the same event time.

\paragraph{Effects: large, sign-correct, unattributable.} The
lifecycle-matched TWFE estimate for the H20 (Table~\ref{tab:main}, panel C)
is $\hat\tau=-2.756$ asinh units (CR1 se $0.652$, $t=-4.23$) with 0 of 6
matched placebos exceeding --- but 6 controls cap the exact test at its
$p=1/7=0.143$ floor. The event study is clean in shape: pre-period
coefficients $+0.19$, $+0.37$, $+0.29$, $+0.01$, $0$; post-period $-0.80$,
$-2.88$, $-4.07$ (Figure~\ref{fig:fig1}c). The naive (unmatched) legs
reproduce the lifecycle failure and calibrate how oversized CR1 is on these
panels: the naive H20 spec has $t=-4.09$ yet 7 of 19 placebo chips produce
larger pseudo-effects ($p=0.400$); naive Oct-23 gives $p=0.727$; naive
components $p=0.203$. Robustness on the matched leg: dropping the 2025Q1
anticipation spike leaves $\hat\tau=-2.799$ ($t=-3.73$) but costs one
placebo exceedance ($p=0.286$); restricting to the 3 non-Chinese controls
gives $\hat\tau=-2.018$ with placebo $p=0.750$ --- the contamination check
fails its placebo test, with the H100/H200's own Blackwell wind-down
($-2.12$) and Trainium2 ($+2.88$) both comparable to the estimate. The
placebo-in-time is clean ($\hat\tau=-0.280$, $t=-0.67$).

\begin{table}[t]
\centering
\scriptsize
\caption{Headline numbers of record, all three legs. Panel A: known-treatment
SuDDDS scans; the declared quarter is 2023Q4 for the Oct-23 panel and 2025Q2
for both H20 panels, and $\hat t_0$ is the scan's maximum-LLR break quarter.
Panel B: unsupervised discovery (no treatment column; $\theta=$ the flow).
Panel C: TWFE effects in asinh units; the exact placebo-in-space $p$ is the
inference of record.}
\label{tab:main}
\setlength{\tabcolsep}{4pt}
\begin{tabular}{lrrrr}
\toprule
\multicolumn{5}{l}{\emph{Panel A: known-treatment scans (date localization)}}\\
 & $\hat t_0$ & $W$ & LLR & scan $p$ \\
\midrule
sales Oct-23 ($7\times 16$; A800+H800), Bernoulli/wcc & 2023Q4 exact & 4 & $4.896$ & $0.030$ \\
\quad normal/ar1\_unit cross-check & 2023Q4 exact & 4 & $4.085$ & $0.14$ \\
sales H20 ($20\times 8$; H20), Bernoulli/wcc & 2025Q2 exact & 4 & $1.532$ & $0.44$ \\
\quad normal/ar1\_unit cross-check & 2025Q2 exact & 2 & $0.914$ & $0.60$ \\
components H20 ($15\times 8$; H20+MI308X), Bernoulli/wcc & 2025Q2 exact & 4 & $2.331$ & $0.13$ \\
\quad normal/ar1\_unit cross-check & 2025Q2 exact & 2 & $1.949$ & $0.23$ \\
\midrule
\multicolumn{5}{l}{\emph{Panel B: unsupervised discovery (top subset; every one a chip generation)}}\\
 & $\hat t_0$ & $W$ & LLR & $p$ \\
\midrule
sales Oct-23: \{A800, H100/H200, H800, TPU v4, TPU v5e\} & 2022Q4 & 4 & $16.16$ & $0.01$ (floor) \\
sales H20: \{H100/H200, H20, MI300A, TPU v5e/v5p, Trainium1\} & 2025Q2 exact & 4 & $17.32$ & $0.05$ \\
components H20: 7-chip old-generation set (incl.\ H20, MI308X) & 2025Q1 & 4 & $16.41$ & $0.02$ \\
\midrule
\multicolumn{5}{l}{\emph{Panel C: TWFE effects (asinh units; exceedances $=$ placebos $\geq|\hat\tau|$)}}\\
 & $\hat\tau$ & CR1 se & $t$ & exceed., $p$ \\
\midrule
matched sales H20 (6 controls alive all 8 q) & $-2.756$ & $0.652$ & $-4.23$ & 0/6, $0.143$ (floor) \\
\quad drop 2025Q1 (anticipation spike) & $-2.799$ & $0.751$ & $-3.73$ & 1/6, $0.286$ \\
\quad non-Chinese controls only (3) & $-2.018$ & $0.939$ & $-2.15$ & 2/3, $0.750$ \\
naive sales Oct-23 (all 5 controls) & $-1.070$ & $1.160$ & $-0.92$ & 7/10, $0.727$ \\
naive sales H20 (all 19 controls) & $-2.347$ & $0.574$ & $-4.09$ & 7/19, $0.400$ \\
naive components H20 (all 13 controls) & $-2.059$ & $0.876$ & $-2.35$ & 15/78, $0.203$ \\
placebo-in-time (pseudo-$t_0$ 2024Q4, pre-only) & $-0.280$ & $0.420$ & $-0.67$ & clean \\
matched Oct-23 analog & \multicolumn{4}{l}{refused: 0 of 5 control chips alive all 16 quarters} \\
\bottomrule
\end{tabular}

\medskip
\parbox{0.95\textwidth}{\footnotesize Notes: all runs \natex{} v0.2.0, seed
0. Scan $p$-values from permutation calibration with $Q=99$ (floor $0.01$);
Bernoulli scans use within-cluster calibration (wcc), normal scans a
dependence-preserving AR(1)-by-unit null. Composition checks pass ($p=1.0$)
in all six panel-A scans; the anticipation test refused (no usable
pre-quarters). CR1 clusters by chip; with $G\leq 7$ the exact
placebo-in-space $p$ is the inference of record, and the naive H20 row
($t=-4.09$, placebo $p=0.400$) calibrates CR1's oversizing on these panels.
Matched controls: Ascend 910B, Ascend 910C, H100/H200, MI300X, Siyuan 590,
Trainium2 (placebo $\hat\tau$: $-1.03$, $+2.06$, $-2.58$, $-1.50$, $+1.63$,
$+1.42$).}
\end{table}

\section{Caveats and Conclusion}\label{sec:conclusion}

Four caveats bound the claims. First, \emph{product-lifecycle churn is a
competing sharp treatment}: unsupervised discovery finds only
chip-generation subsets at the policy quarters --- the Hopper wind-down with
the H20 merged into it at 2025Q2, the Hopper ramp with the A800's birth at
2022Q4 --- so at product granularity the H20 ban and the Blackwell
transition are colinear in event time and cannot be separated inside this
dataset. Second, \emph{control contamination}: 3 of the 6 lifecycle-matched
controls are Chinese accelerators (Ascend 910B/910C, Siyuan 590) that
plausibly \emph{gained} demand from the ban --- a SUTVA violation biasing
$\hat\tau$ negative (their placebo $\hat\tau$'s are $+2.06$, $+1.63$,
$-1.03$) --- while the no-China robustness leg (3 controls) gives
$\hat\tau=-2.02$ with placebo $p=0.75$: the matched leg's 0/6 leans on
contaminated donors. Third, \emph{measurement}: Huawei and Cambricon flows
are annual allocations smeared over quarters --- smooth by construction,
understating placebo dispersion --- and CR1 $t$-statistics are demonstrably
oversized here (the naive H20 leg pairs $t=-4.09$ with placebo $p=0.400$),
which is why the exact placebo-in-space $p$ is the inference of record.
Fourth, \emph{resolution floors}: 6 matched controls cap the placebo test at
$p=1/7=0.143$ and the scan permutation floor is $0.01$ ($Q=99$); no $p$-value
below its floor is claimed. Relatedly, all treated units share one policy
date per round, so scan $p$-values come from the fitted-null panel
calibration rather than cross-unit timing variation --- the Bernoulli
$p=0.030$ for Oct-23 is the defensible headline, with the normal cross-check
($p=0.14$) reported beside it.

The verdict is deliberately asymmetric. Date localization at product
granularity \emph{works}: all six known-treatment scans land on the declared
quarter exactly, the October-2023 round clears its calibrated null
($p=0.030$), and the banned SKUs' shipments hit literal zero within one
quarter --- a policy written against named products is visible in
product-level data on the dot. Attribution of effect \emph{size} at this
granularity does not work, and the pipeline says so rather than
manufacturing a significant number: the H20 collapse is large
($\hat\tau\approx -2.8$ asinh units) and correctly signed, but placebo chips
generate ban-sized effects from ordinary lifecycle churn, the clean-donor
robustness fails its own placebo test, and the matched Oct-23 leg refuses
for want of donors. For an automated natural-experiment pipeline, the
refusals and the correctly identified confound are the point: the machinery
distinguished what these panels can establish --- \emph{when} --- from what
they cannot --- \emph{how much}.

\paragraph{Reproducibility.} All estimates were produced with \natex{}
v0.2.0 \cite{natex} at seed 0 from frozen extracts of the Epoch AI datasets
\cite{epochchipsales,epochchipcomponents} (not committed to the repository).
Figure~\ref{fig:fig1} regenerates deterministically from the committed
\texttt{figures/make\_fig.py}, which asserts the headline scan and effect
estimates against the numbers of record before drawing.

\scriptsize
\begin{thebibliography}{9}
\setlength{\itemsep}{1pt}\setlength{\parskip}{0pt}

\bibitem{allen2022}
Allen, G.~C. (2022).
Choking off China's access to the future of AI.
Center for Strategic and International Studies (CSIS), October 2022.
\url{https://www.csis.org/analysis/choking-chinas-access-future-ai}

\bibitem{bis2023}
Bureau of Industry and Security (2023).
Implementation of additional export controls: Certain advanced computing
items; supercomputer and semiconductor end use; updates and corrections.
Interim final rule, 88~FR~73458, 25 October 2023.

\bibitem{cameron2015}
Cameron, A.~C., and Miller, D.~L. (2015).
A practitioner's guide to cluster-robust inference.
\emph{Journal of Human Resources}, 50(2), 317--372.

\bibitem{chenroth2024}
Chen, J., and Roth, J. (2024).
Logs with zeros? Some problems and solutions.
\emph{Quarterly Journal of Economics}, 139(2), 891--936.

\bibitem{conleytaber2011}
Conley, T.~G., and Taber, C.~R. (2011).
Inference with ``difference in differences'' with a small number of policy
changes.
\emph{Review of Economics and Statistics}, 93(1), 113--125.

\bibitem{crosignani2025}
Crosignani, M., Han, L., Macchiavelli, M., and Silva, A.~F. (2025).
Securing technological leadership? The cost of export controls on firms.
Federal Reserve Bank of New York Staff Report No.~1096 (first circulated
2023 as ``Geopolitical Risk and Decoupling: Evidence from U.S. Export
Controls'').

\bibitem{epochchipcomponents}
Epoch AI (2026).
Data on AI chip components.
\url{https://epoch.ai/data/ai-chip-components} (extract of 2026, quarters
2024Q1--2025Q4).

\bibitem{epochchipsales}
Epoch AI (2026).
Data on AI chip sales.
\url{https://epoch.ai/data/ai-chip-sales} (extract of 2026; estimates
generated January 2026).

\bibitem{fist2023}
Fist, T., and Grunewald, E. (2023).
\emph{Preventing AI Chip Smuggling to China: A Working Paper}.
Center for a New American Security (CNAS), October 2023.

\bibitem{herlands2018}
Herlands, W., McFowland~III, E., Wilson, A.~G., and Neill, D.~B. (2018).
Automated local regression discontinuity design discovery.
In \emph{Proceedings of the 24th ACM SIGKDD International Conference on
Knowledge Discovery and Data Mining (KDD~'18)}.

\bibitem{hillebrandt2026legs}
Hillebrandt, H. (2026).
Bent, bypassed, unbowed: US AI-chip export controls and China's compute.
\natex{} paper collection.
\url{https://haukehillebrandt.github.io/natex/export-controls-three-leg/}

\bibitem{natex}
Hillebrandt, H. (2026).
\emph{natex: automated natural-experiment discovery and estimation} (version
0.2.0). Software.
\url{https://github.com/HaukeHillebrandt/natex}

\bibitem{nvidia2025}
NVIDIA Corporation (2025).
Current report (Form 8-K), 9 April 2025: H20 export-license requirement and
up to \$5.5 billion of associated charges.
US Securities and Exchange Commission.

\end{thebibliography}

\end{document}
